% Part: history
% Chapter: biographies 
% Section: alonzo-church

\documentclass[../../../include/open-logic-section]{subfiles}

\begin{document}

\olfileid[pt]{his}{bio}{chu}

\olsection{Alonzo Church}

\olphoto{church-alonzo}{Alonzo Church}

Alonzo Church nasceu em Washington, DC, em 14 de junho de 1903. Na primeira
infância, um acidente com uma arma de ar comprimido deixou Church cego de um
olho. Ele concluiu o ensino preparatório em Connecticut em 1920 e iniciou sua
formação universitária em Princeton naquele mesmo ano. Concluiu seus estudos
de doutorado em 1927. Após alguns anos no exterior, Church retornou a
Princeton. Church era conhecido por ser extremamente educado e cuidadoso. Sua
escrita no quadro-negro era impecável, e ele preservava artigos importantes
cobrindo-os cuidadosamente com cimento Duco (uma cola transparente). Fora de
suas atividades acadêmicas, gostava de ler revistas de ficção científica e não
tinha receio de escrever aos editores caso identificasse quaisquer imprecisões
nos textos.

As realizações acadêmicas de Church foram notáveis. Junto com seus alunos
Stephen Kleene e Barkley Rosser, ele desenvolveu uma teoria da calculabilidade
efetiva, o cálculo lambda, independentemente do desenvolvimento da máquina de
Turing por Alan Turing. As duas definições de computabilidade são equivalentes
e dão origem ao que hoje é conhecido como a \emph{Tese de Church--Turing}: a de
que uma função dos números naturais é efetivamente computável se e somente se
for computável por meio de uma máquina de Turing (ou do cálculo lambda). Ele
também provou o que hoje é conhecido como o \emph{Teorema de Church}: o
problema da decisão para a validade de fórmulas de primeira ordem é insolúvel.

Church continuou seu trabalho até a velhice. Em 1967, deixou Princeton para ir
à UCLA, onde foi professor até se aposentar, em 1990. Church faleceu em 1º de
agosto de 1995, aos 92 anos.

\begin{reading} 
Para uma breve biografia de Church, veja \citet{EndertonND}. Os escritos
originais de Church sobre o cálculo lambda e o Entscheidungsproblem (Tese de
Church) são \citet{Church1936,Church1936a}.
\citet{Aspray1984} registra uma entrevista com Church sobre a comunidade
matemática de Princeton na década de 1930.
Church escreveu uma série de resenhas de livros para o \emph{Journal of
Symbolic Logic} de 1936 até 1979. Todas estão arquivadas no site de John
MacFarlane \citep{MacFarlane2015}.
\end{reading} 
\end{document}
